\largerpage
\scp{\tsc{Eastern Timor}}{04-03}
We tentatively identify an \tsc{Eastern Timor} node which contains
\ili{Tetun}, \ili{Kairui-Midiki}, Waima{\Q}a, and \ili{Naueti}.
The exact relationship of these languages with regards to one
another and to \tsc{\ili{Lakalei}-Idate} is not completely clear,
and more work may lead us to revise this classification.
\tsc{Eastern Timor} is defined by three sound changes: *R > {\0};
*j > \it{r}; and *ŋ > \it{n}.
Each of these sound changes is also found in \tsc{\ili{Lakalei}-Idate}
(\srf{04-04}), though not to the same extent or in the same words as \tsc{Eastern Timor}.
Nonetheless, at a higher level it may be possible to place \tsc{Eastern
Timor} and \tsc{\ili{Lakalei}-Idate} together.
%in a higher node within \tsc{Timor-Babar}.
Because of this possibility,
\ili{Lakalei} and \ili{Idate} cognates are given throughout this section for comparison.

\largerpage
One difficulty faced in unravelling the history of the languages
of central and eastern Timor is the role played by \ili{Tetun} [tet].
\ili{Tetun} is socially dominant in central
and eastern Timor and has functioned as a regional lingua franca
for centuries.\footnote{
		In the Belu regency of Indonesian Timor,
		one \ili{Welaun} (\crf{ch:CT}) speaker explained to Edwards that just
		as Indonesian is the unifying language (\it{bahasa persekutuan})
		of the country, so \ili{Tetun} is the unifying language of Belu regency.}
This function
is continued in Timor-Leste as one of the national languages
is \ili{Tetun} Dili [tdt] -- a \ili{Tetun}-based creole.
Many terms in languages of this region are identical to \ili{Tetun}
cognates, and it is not always clear whether a particular term is due to
parallel development and/or borrowing from \ili{Tetun}.

The first change shared by \tsc{Eastern Timor} languages is *j
> \it{r} in most words.
Examples are given in \trf{04-09} along with cognates in \ili{Lakalei}
and \ili{Idate} which have *j > \it{l}, \it{r}.
Reflexes of *huaji ‘younger sibling’ are irregular
and have *j > \it{l} in \ili{Tetun}
and \ili{Habun}; \ili{Tetun} \it{ali {\tl} alik {\tl} alin},
\ili{Habun} \it{walin}, \ili{Midiki} \it{warin},
Waima{\Q}a \it{wari}, \ili{Naueti} \it{ware}.
\ili{Tetun} \it{ali {\tl} alik {\tl} alin} ‘younger sibling’ is also
irregular in showing initial *hua > **wa > {\0}.

\begin{table}\tcp{\tsc{Eastern Timor} and \tsc{\ili{Lakalei}-Idate} *j}{04-09}
\begin{adjustbox}{width=\textwidth}
	\begin{threeparttable}\st{2pt}
\begin{tabular}{llllllll}\lsptoprule
local gl.	&	sun	&	how m.	&	nose	&	dry	&	foam	&	bile	&	name\\
PMP	&	*qalə\tbr{j}aw	&	*pi\tbr{j}a	&	*i\tbr{j}uŋ	&	*ma\tbr{j}a	&	*bu\tbr{j}əq	&	*qapə\tbr{j}u	&	*ŋa\tbr{j}an\\ \midrule
E. \ili{Tetun}	&	{\hr}lo\tbr{r}o	&	{\hr}hi\tbr{r}a	&	{\hr}inu\tbr{r}\su{†}	&	{\hr}ma\tbr{r}an	&	{\hr}fu\tbr{r}in	&	{\hr}naʔan ho(u)\tbr{r}un	&	{\hr}na\tbr{r}an\\
\ili{Habun}	&	{\hr}lo\tbr{r}on	&	{\hr}hi\tbr{r}a	&	{\hr}i\tbr{r}un	&	{\hr}ma\tbr{r}an	&		&		&	{\hr}na\tbr{r}an\\
\ili{Midiki}	&	{\hr}la\tbr{r}a	&	{\hr}kai-hi\tbr{r}a	&	\cg{(tukai)}	&	{\hr}ma\tbr{r}i	&		&		&	\cg{(nai-n)}\\
Waima{\Q}a	&	{\hr}ʔla\tbr{r}a	&	{\hr}hi\tbr{r}e	&	\cg{(turu-kai)}	&	{\hr}hma\tbr{r}e	&	{\hr}wu\tbr{r}o	&	{\hr}he\tbr{r}u	&	\cg{(nai)}\\
\ili{Naueti}	&	{\hr}ʔla\tbr{r}a	&	{\hr}kai-hi\tbr{r}a	&	{\hr}i\tbr{r}u, i\tbr{r}ukai	&	{\hr}ma\tbr{r}e	&	{\hr}wu\tbr{r}an	&	{\hr}ye\tbr{r}u	&	\cg{(nai)}\\ \midrule
\ili{Lakalei}	&	{\hr}lelo	&	{\hr}ila	&	{\hr}inun	&	\cg{(kiʔas)}	&		&		&	{\hr}na\tbr{r}an\\
\ili{Idate}	&	{\hr}lelo	&	{\hr}yila	&	{\hr}yinun	&	{\hr}ma\tbr{r}an	&	{\hr}hu\tbr{r}an	&	{\hr}melu-k\su{‡}	&	{\hr}na\tbr{r}an\\
		\lspbottomrule
	\end{tabular}
		\begin{tablenotes}\tnsize
			\item[†]	{\ili{Tetun} \it{inur} ‘nose’ has consonant metathesis.}
			\item[‡]	{\ili{Idate} \it{meluk} ‘bitter’ is from intermediate **ma-pəju.
								Compare East \ili{Tetun} \it{moruk} ‘bitter’.}
		\end{tablenotes}
	\end{threeparttable}
\end{adjustbox}
\end{table}

The \ili{Midiki}, Waima{\Q}a, and \ili{Naueti} terms for ‘name’ in \trf{04-09} are
probably borrowed from neighbouring non-Austronesian \ili{Makasae}
and are not directly inherited from PMP *ŋajan.
Proto-Timor *-nej ‘name’ has been reconstructed for the Papuan
languages of Timor \citep[117]{SchapperEngelenhoven2017}.\footnote{
		\ili{Makasae} \it{nai} and/or Proto-Timor *-nej
		may themselves be borrowed from an Austronesian source.
		The most similar words regionally are from languages of Atauro island (\srf{04-05}); e.g.
		\ili{Raklungu} \it{ŋai}, \ili{Hresuk} \it{ŋae},
		Dadu{\Q}a \it{nae} ‘name’.} 

Word initially and medially PMP *R is usually lost in \tsc{Eastern Timor} languages.
Examples are given in \trf{04-10}.
\ili{Lakalei} and \ili{Idate} have *R > \it{r,} {\0}.
\ili{Habun} has *R > \it{ʔ} in *diRus > \it{riʔu} ‘bathe’
and this may be intermediate between *R
and {\0}, or a later insertion.

\begin{table}\tcp{\tsc{Eastern Timor} and \tsc{\ili{Lakalei}-Idate} *R}{04-10}
%\begin{adjustbox}{width=\textwidth}
	\begin{threeparttable}\st{2.2pt}
	\begin{tabular}{lllllllll}\lsptoprule
local gl.	&	blood	&	bathe	&	stand	&	land	&	bone	&	red	&	new	&	house\\
PMP	&	*da\tbr{R}aq	&	*di\tbr{R}us	&	*di\tbr{R}i	&	*da\tbr{R}əq	&	*du\tbr{R}i	&	*ma-i\tbr{R}aq	&	*baqə\tbr{R}u	&	*\tbr{R}umaq\\ \midrule
E. \ili{Tetun}	&	{\hr}raa-n	&	{\hr}ha-riis	&	{\hr}ha-rii	&	{\hr}rai	&	{\hr}rui-n	&	{\hr}mean	&	{\hr}foun	&	{\hr}uma\\
\ili{Habun}	&	{\hr}raa-n	&	{\hr}n-riʔu	&	{\hr}na-rii\su{†}	&	{\hr}rai	&	{\hr}rui-n	&	{\hr}mean	&	{\hr}woun	&	\\
\ili{Midiki}	&	{\hr}raa-n	&	{\hr}riu	&	{\hr}hnii	&	{\hr}raa	&	{\hr}rui-n	&	{\hr}mee	&	\cg{(morin)}	&	\\
Waima{\Q}a	&	{\hr}raa	&	{\hr}riu	&	{\hr}hnii	&	{\hr}ria, rie	&	{\hr}rui	&	{\hr}mee	&	\cg{(mori)}	&	{\hr}umo\\
\ili{Naueti}	&	{\hr}raa	&	{\hr}riu	&	\cg{(ditu)}	&	{\hr}rea	&	{\hr}rui	&	{\hr}mee	&	\cg{(mori)}	&	{\hr}uma\\ \midrule
\ili{Lakalei}	&	{\hr}raa-n	&	{\hr}rius	&	{\hr}amiri	&	{\hr}lare	&	{\hr}rui-n	&	{\hr}meren	&	\cg{(morin)}	&	\cg{(lehu)}\\
\ili{Idate}	&	{\hr}raa-n	&	{\hr}a-rius	&	{\hr}ba-rii	&	{\hr}larek	&	{\hr}lurik	&	{\hr}meran	&	\cg{(morin)}	&	\cg{(ada)}\\
		\lspbottomrule
	\end{tabular}
		\begin{tablenotes}\tnsize
			\item[†] {\ili{Habun} \it{na-rii} is ‘make stand, build’.}
		\end{tablenotes}
	\end{threeparttable}
%\end{adjustbox}
\end{table}

\begin{table}\tcp{\tsc{Eastern Timor} and \tsc{\ili{Lakalei}-Idate} word final *R}{04-11}
%\begin{adjustbox}{width=\textwidth}
	\begin{threeparttable}\st{5.5pt}
	\begin{tabular}{lllllll}\lsptoprule
local gloss	&	egg	&	tail	&	bamboo	&	hungry	&	fry	&	coconut\\
PMP		&	*qatəl\tbr{uR}	&	*ik\tbr{uR}	&	*qa\tbr{uR}	&	*lap\tbr{aR}	&	*saŋəl\tbr{aR}\su{†}	&	*ni\tbr{uR}\\	\midrule
East \ili{Tetun}	&	{\hqa}tolun	&	{\hr}ikun	&	{\hq}au	&	{\hr}hamlaha	&	{\hr}sona	&	{\hr}nuu\\
\ili{Habun}		&	{\hqa}tel\tbr{i}n&	{\hr}ik\tbr{i}n&	{\hq}a\tbr{i}	&	{\hr}na-nlah\tbr{e}	&	{\hr}sin\tbr{i}	&	{\hr}no\tbr{e}\\
\ili{Midiki}		&	{\hqa}htelo	&	{\hr}iku	&	{\hq}aa	&	{\hr}laha	&	{\hr}sena	&	{\hr}neo\\
Waima{\Q}a	&	{\hqa}htelu	&	{\hr}iku	&	{\hq}a\tbr{e}	&	{\hr}laha	&	{\hr}sena	&	{\hr}ne\tbr{e}\\
\ili{Naueti}		&	{\hqa}htelu	&	{\hr}iku	&	{\hq}a\tbr{e}	&	\cg{(luba)}	&	{\hr}sena	&	{\hr}ne\tbr{e}\\ \midrule
\ili{Lakalei}	&	{\hqa}telor	&	{\hr}hiʔon	&			&			&		&	\\
\ili{Idate}		&	{\hqa}telor	&	{\hr}yiʔan	&	{\hq}oor	&	\cg{(beli)}	&		&	{\hr}noor\\
		\lspbottomrule
	\end{tabular}
		\begin{tablenotes}\tnsize
			\item[†]	{Reflexes of *saŋəlaR are via intermediate **sənaR.
								\ili{Habun} \it{sini} is ‘burn’ and may instead
								be connected with \ili{Tetun} Dili \it{sunu} ‘burn’.}
		\end{tablenotes}
	\end{threeparttable}
%\end{adjustbox}
\end{table}

\largerpage
There is evidence that *R > {\0} went through an intermediate
glide stage **y in at l\ili{east} some words for most of these languages.
Firstly, the vowels of all reflexes of *daRəq ‘land’ are not
fully explained by development from intermediate **daə(q) ---
with the exception of \ili{Midiki} \it{raa} ‘land’.
In particular, the Waima{\Q}a
and \ili{Naueti} words with a penultimate front vowel
probably point to a form with a medial glide **y.
This may be due to conflation with PMP *daya
`upriver, toward the interior', that is `landward'.

Secondly, word-final *R often triggers fronting of the previous vowel in \ili{Habun}.
Examples are shown in \trf{04-11}.
There are three examples of *uR{\#} > \it{i},
one of *aR{\#} > \it{e}
and one which probably shows *aR > \it{i}.
Waima{\Q}a and \ili{Naueti} also show otherwise unexplained fronting
of the final vowel in *qauR > \it{ae} ‘bamboo’.
These examples point to word-final *R > **y > {\0} in \ili{Habun}, Waima{\Q}a, and \ili{Naueti}.

\begin{table}\tcp{\tsc{Eastern Timor} and \tsc{\ili{Lakalei}-Idate} *ŋ > \it{n}}{04-12}
\begin{adjustbox}{width=\textwidth}
	\begin{threeparttable}\st{2pt}
	\begin{tabular}{lllllll}\lsptoprule
local gloss	&	palate	&	young	&	ear	&	fry	&	swim	&	sky\\
PMP	&	*\tbr{ŋ}adas\su{†}	&	*\tbr{ŋ}uda	&	*tali\tbr{ŋ}a	&	*sa\tbr{ŋ}əlaR	&	*na\tbr{ŋ}uy	&	*la\tbr{ŋ}it\su{¶}\\ \midrule
E. \ili{Tetun}	&	{\hr}kna{\tl}\tbr{n}arak	&	{\hr}\tbr{n}urak	&	\cg{(tilun)}	&	{\hr}so\tbr{n}a	&	{\hr}na\tbr{n}i	&	\cg{(lalehan)}\\
\ili{Fehan} \ili{Tetun}	&	{\hr}ka\tbr{n}arak	&	{\hr}\tbr{n}urak	&	\cg{(tilun)}	&	{\hr}so\tbr{n}a	&	{\hr}na\tbr{n}i	&	{\hr}\tbr{n}ali\\
\ili{Habun}	&		&		&	{\hr}hi\tbr{n}a-n	&	{\hr}si\tbr{n}i	&	{\hr}na\tbr{n}i	&	\cg{(lalehan)}\\
\ili{Midiki}	&		&	{\hr}\tbr{n}ura	&	\cg{(nihura)}	&	{\hr}se\tbr{n}a	&	{\hr}na\tbr{n}i	&	{\hr}\tbr{n}ali\\
Waima{\Q}a	&	{\hr}ʔ\tbr{n}ire	&	{\hr}\tbr{n}uro	&	\cg{(ʔnihu)}	&	{\hr}se\tbr{n}a	&	\cg{(klee)}	&	{\hr}\tbr{n}ali\\
\ili{Naueti}	&	{\hr}\tbr{n}ira	&	{\hr}\tbr{n}ura	&	\cg{(gara)}	&	{\hr}se\tbr{n}a	&	\cg{(lele)}	&	{\hr}\tbr{n}ali\\ \midrule
\ili{Lakalei}	&		&	\cg{(losan)}	&	{\hr}ti\tbr{n}a roo	&	{\hr}se\tbr{n}e	&	{\hr}na\tbr{n}i	&	\cg{(lalehan)}\\
\ili{Idate}	&		&	\cg{(kaek)}	&	{\hr}tali\tbr{n}a	&	{\hr}se\tbr{n}ar	&	{\hr}na\tbr{n}i	&	\cg{(lalehan)}\\ \midrule \midrule
**gloss	&	{\hr}gills, palate	&	{\hr}young	&	{\hr}ear	&	{\hr}fry	&	{\hr}swim	&	{\hr}heaven\\
PR-M	&	**\tbb{ŋg}adas	&	**\tbb{ŋ}ura-k	&	**ta\tbb{ŋ}ila\su{Δ}	&	**se\tbb{ŋ}a	&	**na\tbb{ŋ}e	&	**la\tbr{n}i\\ \midrule
\ili{Termanu}	&		&	{\hr\hr}\tbr{n}ula-k	&	{\hr\hr}\tbb{ŋg}ela-k	&	{\hr\hr}se{\tl}se\tbr{n}a	&	{\hr\hr}a\tbb{ŋ}e	&	{\hr\hr}la{\tl}lai\\
Ba{\Q}a	&	{\hr\hr}ŋga{\tl}\tbb{ŋg}alas	&	{\hr\hr}\tbr{n}ula-k	&	{\hr\hr}\tbb{ŋg}ela-k	&	{\hr\hr}se{\tl}se\tbr{n}a	&	{\hr\hr}a\tbr{n}e	&	{\hr\hr}la{\tl}lai\\
\ili{Dela}-O.	&	{\hr\hr}\tbb{ŋg}ara-ʔ	&	{\hr\hr}\tbr{n}ure-ʔ	&	{\hr\hr}\tbb{ŋg}ela-ʔ	&	{\hr\hr}se{\tl}se\tbr{n}a	&	{\hr\hr}na\tbr{n}e	&	{\hr\hr}la\tbr{n}i\\
\ili{Kotos} 	&	{\hr\hr}ʔ\tbb{k}are-f	&	{\hr\hr}mainu\tbb{k}eʔ\su{‡}	&	{\hr\hr}na-t\tbr{n}ina	&		&	\hr\cg{(na-bhaʔe)}	&	\hr\cg{(neno)}\\
Kusa-M.	 	&		&	{\hr\hr}ma\tbb{k}ukaʔ\su{‡}	&		&	{\hr\hr}ta-see\tbb{k}\su{ǁ}	&	\hr\cg{(roeʔ)}	&	\hr\cg{(neno)}\\
		\lspbottomrule
	\end{tabular}
		\begin{tablenotes}\tnsize
			\item[†]	{\ili{Fehan} \ili{Tetun} \it{kanarak} ‘throat’, \ili{Naueti} \it{nira} ‘gums’.
								\ili{Naueti} and Waima{\Q}a have irregular penultimate \it{i}.}
			\item[‡]	{\ili{Kotos} \ili{Amarasi} \it{mainukeʔ} has consonant metathesis from intermediate **ma-nuŋe-k.
								\ili{Kusa-Manea} \it{makukaʔ} is probably from intermediate
								**ma-ŋuŋa-k with assimilation of *n > **ŋ.}
			\item[Δ]	{\citet[302]{Edwards2021} reconstructs **ŋgela-k `earwax'
									(which is the meaning found in the Rote languages).
									However, identification of \ili{Kotos} \ili{Amarasi} \it{na-tnina} `listen'
									as connected \citep[141]{Tan2023},
									points to Proto-Rote-\ili{Meto} **taŋila.
									\citet{BlustTrussel2020} reconstruct doublets
									*taliŋa and *taŋila.}
			\item[{ǁ}]	{\ili{Kusa-Manea} \it{ta-seek} ‘fry’ has final CV metathesis
									from (currently unattested) *[-seke].}
			\item[¶]	{Proto-Eastern Timor **nali has consonant metathesis from earlier **lani.
								\ili{Fehan} \ili{Tetun} \it{nali} is given as ‘sky in the cosmological sense’.
								(The normal word for ‘sky’ is \it{laleʔan}.) Rote reflexes of
								*laŋit are given with the meaning ‘heaven, firmament, air space’.
								\ili{Termanu} and Ba{\Q}a have irregular *ŋ > {\0}.}
		\end{tablenotes}
	\end{threeparttable}
\end{adjustbox}
\end{table}

\newpage
Supporting evidence for \tsc{Eastern Timor} comes from *ŋ > \it{n}.
\ili{As} discussed at the beginning of this chapter,
this change is widespread throughout the \tsc{Timor-Babar} languages
and probably spread by diffusion.
Nonetheless, the \tsc{Eastern Timor} languages appear to have
completed this change while other nodes have not.
Thus, *ŋ > **n can be posited at the level of Proto-Eastern Timor.
 Examples of *ŋ > \it{n} in \tsc{Eastern Timor} where \tsc{Rote-Meto}
(\srf{04-02}) usually has *ŋ > **ŋ,
**ŋg are given in \trf{04-12} to show the contrast.
(In the case of *laŋit ‘sky’,
\ili{Teun}, \ili{Nila}, and \ili{Serua} have \it{lakti} ‘sky’.
Comparison of this word with others, such as \ili{Luang} \it{lyanti}
points to Proto-Southwest Maluku **ŋ for this word.) 

\ili{Midiki}, Waima{\Q}a and \ili{Naueti} can be placed together based on
*-ay/*-aw > \it{a} and *-uy > \it{u} in most words.
\ili{Tetun} has *-ay > \it{e},
*-aw > \it{o} and *-uy > \it{i}.
\ili{Habun} also has *-aw > \it{a} in two words,
and *-aw > \it{o} in *qaləjaw > \it{loron} ‘sun’ (probable loan from \ili{Tetun}).
\ili{Habun} has *-ay > \it{e} in all three known reflexes,
but all these words are identical to \ili{Tetun}
and could be loans.
\ili{Habun} has *-uy > \it{i} in one word which is not identical to
\ili{Tetun}, and in this regard these languages pattern together.
Examples of final vowel-glide sequences in \tsc{Eastern Timor}
are shown in \trf{04-13}.
\ili{Tetun}, \ili{Habun}, \ili{Midiki} all have *uy > \it{i} in *naŋuy > \it{nani}
‘swim’ (see \trf{04-12}).

\begin{table}\tcp{\tsc{Eastern Timor} and \tsc{\ili{Lakalei}-Idate} final VG sequences}{04-13}
\begin{adjustbox}{width=\textwidth}
	\st{2pt}\begin{tabular}{lllllllll}\lsptoprule
local gl.		&	die	&	liver	&	go up	&	rattan	&	mouse	&	walk	&	day, sun	&	pig\\
PMP					&	*mat\tbr{ay}		&	*qat\tbr{ay}					&	*sak\tbr{ay}		&	*qu\tbr{ay}		&	*balab\tbr{aw}			&	*lak\tbr{aw}			&	*qaləj\tbr{aw}		&	*bab\tbr{uy}\\ \midrule
E. \ili{Tetun}		&	{\hr}mat\tbr{e}	&	{\hq}at\tbr{e}n	&	{\hr}saʔ\tbr{e}	&	{\hq}o\tbr{e}	&	{\hr}lah\tbr{o}			&	{\hr}laʔ\tbr{o}		&	{\hr}lor\tbr{o}n	&	{\hr}fah\tbr{i}\\
\ili{Habun}				&	{\hr}mat\tbr{e}	&	{\hq}at\tbr{e}n	&									&	{\hq}o\tbr{e}	&	{\hr}lab\tbr{a}			&	{\hr}n-lak\tbr{a}	&	{\hr}lor\tbr{o}n	&	{\hr}wab\tbr{i}\\
\ili{Midiki}			&	{\hr}hmat\tbr{a}&	{\hq}at\tbr{i}n	&									&	{\hq}u\tbr{a}	&	{\hr}hlow\tbr{a}		&	{\hr}hlak\tbr{a}	&	{\hr}lar\tbr{a}		&	\\
Waima{\Q}a	&	{\hr}mat\tbr{a}	&	{\hq}at\tbr{a}	&	{\hr}sak\tbr{a}	&	{\hq}u\tbr{o}	&	{\hr}lu\tbr{a}			&	{\hr}lak\tbr{a}		&	{\hr}ʔlar\tbr{a}	&	{\hr}wa\tbr{u}\\
\ili{Naueti}			&	{\hr}mat\tbr{a}	&	{\hq}at\tbr{a}	&	{\hr}sak\tbr{a}	&	{\hq}u\tbr{a}-kai	&	{\hr}lo\tbr{a}	&	{\hr}lak\tbr{a}		&	{\hr}ʔlar\tbr{a}	&	{\hr}wo\tbr{u}\\ \midrule
\ili{Lakalei}			&	{\hr}mate				&	{\hq}aten				&									&								&											&										&	{\hr}lelo					&	\\
\ili{Idate}				&	{\hr}mate				&	{\hq}ater				&	{\hr}saen				&	{\hr}ua				&	{\hr}malaho					&	{\hr}lau					&	{\hr}lelo					&	{\hr}hahi\\
		\lspbottomrule
	\end{tabular}
\end{adjustbox}
\end{table}

Additional evidence for placing \ili{Midiki},
Waima{\Q}a, and \ili{Naueti} together comes from shared penultimate
*ə > \it{e} in most words.
\ili{Tetun} has penultimate *ə > \it{o}.
While \ili{Habun} has penultimate *ə > \it{o} before back vowels
and penultimate *ə > \it{e} before other vowels.
The reflexes of *ə in final syllables are more complex.
\ili{Tetun} has *ə > \it{a}
/σ\und{σ}{\#} in most cases,
while \ili{Habun} has a split of *ə > \it{e,
a} /σ\und{σ}{\#}, perhaps conditioned by the quality of the previous vowel.
Waima{\Q}a and \ili{Naueti} have *ə > \it{e} /σ\und{σ}{\#} in most
cases, while \ili{Midiki} has *ə > \it{i, a}.
Examples of *ə in \tsc{Eastern Timor} are given in \trf{04-14}.

\begin{table}\tcp{\tsc{Eastern Timor} and \tsc{\ili{Lakalei}-Idate} *ə}{04-14}
%\begin{adjustbox}{width=\textwidth}
	\st{1.75pt}\begin{tabular}{llll@{ }lll@{ }l}\lsptoprule
local gl.	&	egg	&	hear	&	three	&	new	&	full	&	itch	&	inside\\
PMP				&	*qat\tbr{ə}luR		&	*d\tbr{ə}ŋəR	&	*t\tbr{ə}lu	&	*baq\tbr{ə}Ru	&	*(ma)-p\tbr{ə}nuq	&	*(ma)-gat\tbr{ə}l	&	*dal\tbr{ə}m\\	\midrule
E. \ili{Tetun}		&	{\hqa}t\tbr{o}lun	&	{\hr}r\tbr{o}na	&	{\hr}t\tbr{o}lu	&	{\hr}f\tbr{o}un	&	{\hr}k\tbr{o}nu(n)	&	{\hr}kat\tbr{a}l, kat\tbr{a}r	&	{\hr}lar\tbr{a}n\\
\ili{Habun}			&	{\hqa}t\tbr{e}lin	&	{\hr}n-r\tbr{e}ne	&	{\hr}t\tbr{o}lu	&	{\hr}w\tbr{o}un	&	{\hr}h\tbr{o}nu	&	{\hr}hat\tbr{a}	&	{\hr}rel\tbr{e}n\\
\ili{Midiki}		&	{\hq}ht\tbr{e}lo	&	{\hr}d\tbr{i}na	&	{\hr}kai-t\tbr{e}lu	&	\cg{(morin)}	&	{\hr}b\tbr{e}nu	&	{\hr}mat\tbr{i}	&	{\hr}lal\tbr{i}n\\
Waima{\Q}a&	{\hq}ht\tbr{e}lu	&	{\hr}d\tbr{e}ne	&	{\hr}t\tbr{e}lu	&	\cg{(hmori)}	&	{\hr}b\tbr{e}no	&	{\hr}ʔmat\tbr{e}	&	{\hr}lal\tbr{e}\\
\ili{Naueti}		&	{\hq}ht\tbr{e}lu	&	{\hr}d\tbr{e}ne	&	{\hr}kai-t\tbr{e}lu	&	\cg{(mori)}	&	{\hr}b\tbr{e}nu	&	{\hr}ʔnat\tbr{e}	&	{\hr}lal\tbr{e}\\	\midrule
\ili{Lakalei}		&	{\hqa}t\tbr{e}lor	&	\cg{(lii)}	&	{\hr}t\tbr{e}lu	&	\cg{(morin)}	&	{\hr}b\tbr{e}nu	&		&	{\hr}lal\tbr{a}n\\
\ili{Idate}			&	{\hqa}t\tbr{e}lor	&	\cg{(sei)}	&	{\hr}t\tbr{e}lu	&	\cg{(mori)}	&	{\hr}na-b\tbr{e}no	&	\cg{(kalatak)}	&	{\hr}lal\tbr{a}n\\
		\lspbottomrule
	\end{tabular}
%\end{adjustbox}
\end{table}

Another sound change shared between \ili{Habun},
\ili{Midiki}, Waima{\Q}a, and \ili{Naueti} is initial *ma-p > \it{b}.
This sound change is common in the region
and also occurs in \tsc{\ili{Lakalei}-Idate} (\srf{04-04}),
\ili{Galolen} (\srf{04-05}), and \tsc{Nuclear \ili{Central} Timor} (\crf{ch:CT}).
\ili{Midiki}, Waima{\Q}a, and \ili{Naueti} also share medial *q > \it{ʔ}.
\ili{Habun} appears to have medial *q > \it{h},
though only one putative example has been identified.
\ili{Tetun} has *ma-p > \it{m}
and *q > {\0}.
Examples of PMP *ma-p
and *q are given in \trf{04-15}.

\begin{table}\tcp{\tsc{Eastern Timor} and \tsc{\ili{Lakalei}-Idate} *ma-p and *q}{04-15}
%\begin{adjustbox}{width=\textwidth}
	\st{2pt}\begin{tabular}{llll|llll}\lsptoprule
local gl.	&	white	&	hot	&	bitter	&	faeces	&	language	&	thigh	&	trunk\\
PMP	&	*\tbr{ma-p}utiq	&	*\tbr{ma-p}anas	&	*\tbr{ma-p}a\tbr{q}it	&	*ta\tbr{q}i	&	*li\tbr{q}əR	&	*pa\tbr{q}a	&	*pu\tbr{q}un\\ \midrule
E. \ili{Tetun}	&	{\hr}\tbr{m}utin	&	{\hr}\tbr{m}anas	&	\cg{(moruk)}	&	{\hr}teen	&	{\hr}lia	&	\cg{(kelen)}	&	{\hr}huun\\
\ili{Habun}	&	{\hr}\tbr{b}utin	&	\cg{(hutun)}	&		&	{\hr}ta\tbr{h}ek	&		&	\cg{(kelen)}	&	\\
\ili{Midiki}	&	{\hr}\tbr{b}uti	&	\cg{(mutu)}	&		&	{\hr}ta\tbr{ʔ}i	&		&	{\hr}ha\tbr{ʔ}a	&	{\hr}kai ho\tbr{ʔ}o\\
Waima{\Q}a	&	{\hr}\tbr{b}uto	&	\cg{(ʔmutu)}	&	{\hr}\tbr{b}a\tbr{ʔ}i	&	{\hr}ta\tbr{ʔ}i	&	{\hr}li\tbr{ʔ}a	&	{\hr}ha\tbr{ʔ}a	&	{\hr}ho\tbr{ʔ}o\\
\ili{Naueti}	&	{\hr}\tbr{b}uta	&	\cg{(ʔmutu)}	&	{\hr}\tbr{b}a\tbr{ʔ}i	&	{\hr}ta\tbr{ʔ}i	&	{\hr}le\tbr{ʔ}a, lea	&		&	{\hr}kai ho\tbr{ʔ}o\\ \midrule
\ili{Lakalei}	&	{\hr}\tbr{b}utin	&	{\hr}\tbr{b}anas	&		&		&		&		&	{\hr}ai uun\\
\ili{Idate}	&	{\hr}\tbr{b}utin	&	{\hr}\tbr{b}anas	&	\cg{(meluk)}	&	{\hr}ta\tbr{h}ek	&		&		&	\\
		\lspbottomrule
	\end{tabular}
%\end{adjustbox}
\end{table}

\ili{Habun} and \ili{Midiki} share initial *z > \it{y}
and medial *z > {\0}.
Other Eastern Timor languages have *z > \it{d}.
\ili{Habun}, \ili{Midiki}, Waima{\Q}a, and \ili{Naueti} share initial *b > \it{w},
with \ili{Midiki} having *b > **w > \it{w, {\0}} /{\gap}u.
Similarly, word medially Waima{\Q}a, \ili{Naueti},
and \ili{Midiki} share *b > **w > \it{w,} {\0} along with rounding
of the previous vowel.
\ili{Habun} retains *b = \it{b} word medially,
though this may be due to a later change of medial **w > \it{b},
thus *b > **w > \it{b}.
\ili{Tetun} has initial *b > \it{f}
and medial *b > \it{h}.
Examples of the reflexes of *z
and *b in \tsc{Eastern Timor} are given in \trf{04-16}.
(See also reflexes of *balabaw ‘mouse’ in \trf{04-13}.)

\begin{table}\tcp{\tsc{Eastern Timor} and \tsc{\ili{Lakalei}-Idate} *z and *b}{04-16}
%\begin{adjustbox}{width=\textwidth}
	\begin{threeparttable}\st{2pt}
	\begin{tabular}{ll@{}ll|lllll}\lsptoprule
local gl.	&	road	&	far	&	rain	&	stone	&	moon	&	fruit	&	pig	&	mouth\\
PMP	&	*\tbr{z}alan	&	*\tbr{z}auq	&	*qu\tbr{z}an	&	*\tbr{b}atu	&	*\tbr{b}ulan	&	*\tbr{b}uaq	&	*\tbr{b}a\tbr{b}uy	&	*\tbr{b}aq\tbr{b}aq\\ \midrule
E. \ili{Tetun}	&	{\hr}\tbr{d}alan	&	{\hr}\tbr{d}ook	&	{\hr}u\tbr{d}an	&	{\hr}\tbr{f}atu\cg{(k)}	&	{\hr}\tbr{f}ulan	&	{\hr}\tbr{f}uan	&	{\hr}\tbr{f}a\tbr{h}i	&	\cg{(ibun)}\\
\ili{Habun}	&	{\hr}\tbr{y}ala	&	\cg{(naun)}	&	{\hr}oe	&	{\hr}\tbr{w}atu	&	{\hr}\tbr{w}ulan	&	{\hr}\tbr{w}uan	&	{\hr}\tbr{w}a\tbr{b}i	&	{\hr}\tbr{w}a\tbr{b}an\\
\ili{Midiki}	&	{\hr}\tbr{y}ala	&	{\hr}\tbr{y}au	&	{\hr}ua	&	{\hr}\tbr{w}atu	&	{\hr}\tbr{w}ula	&	{\hr}ua	&		&	\cg{(noe)}\\
Waima{\Q}a	&	{\hr}\tbr{d}ala	&	{\hr}\tbr{d}au	&	{\hr}u\tbr{d}o	&	{\hr}\tbr{w}atu	&	{\hr}\tbr{w}ulo	&	{\hr}\tbr{w}uo	&	{\hr}\tbr{w}au	&	{\hr}\tbr{w}uwa\\
\ili{Naueti}	&	{\hr}\tbr{d}ala	&	{\hr}\tbr{d}au	&	{\hr}u\tbr{d}a	&	{\hr}\tbr{w}atu	&	{\hr}\tbr{w}ula	&	{\hr}\tbr{w}ua	&	{\hr}\tbr{w}ou	&	\cg{(nunu)}\\ \midrule
\ili{Lakalei}	&	{\hr}lal usu	&	{\hr}roon	&	{\hr}huʔan	&	{\hr}hatuk	&	{\hr}hulkai	&	{\hr}huan	&		&	\cg{(ibar)}\\
\ili{Idate}	&	{\hr}lala-matak\su{†}	&	{\hr}arook	&	{\hr}wari	&	{\hr}hatu	&	{\hr}hula	&	{\hr}huan	&	{\hr}hahi	&	\cg{(yibar)}\\
		\lspbottomrule
	\end{tabular}
		\begin{tablenotes}\tnsize
			\item[†] {\ili{Idate} \it{lala-matak} is ‘door’.}
		\end{tablenotes}
	\end{threeparttable}
%\end{adjustbox}
\end{table}

